% !TeX root = ../深度学习与大语言模型技术原理.tex
% 由主文档通过 \input 载入；请编译主文档。
\setcounter{chapter}{16} % 保留原书章号。
\chapter{智能体循环、记忆与规划}
\label{ch:agents}

\section{状态、动作与终止条件}

智能体（agent）围绕目标，根据观察选择动作，并利用环境反馈更新后续决策。工作流（workflow）可以预先固定执行图，智能体则把部分路径选择交给模型；实际系统常组合二者。一个可控循环维护任务状态、历史观察、已完成子目标、未解决问题和剩余预算，在每次动作后确认状态变化，再选择继续、恢复或结束。

将系统写成状态机有助于明确语义：待决策、待执行、待观察、待恢复和已终止具有不同允许操作。终止条件包括独立验证成功、明确不可继续和预算耗尽。最后一种状态不能等同于成功，模型说“完成”也不能替代对真实环境的检查。

长任务应在安全边界保存检查点，把已提交副作用和可重放的纯计算分开记录。循环检测可以比较规范化动作、相同观察与状态摘要；连续重复且没有进展时，应改变策略或停止，而非无限增加提示中的自我提醒。

\section{记忆的类型与更新规则}

会话记忆（conversation memory）保存交互历史，工作记忆（working memory）保存当前任务状态，情景记忆（episodic memory）记录过去事件，语义记忆（semantic memory）保存可复用知识。这是功能划分，并不意味着必须建立四个数据库。实现可以使用摘要、向量索引、键值存储和结构化表。

记忆写入比读取更需要谨慎，因为一次错误可能影响后续很多任务。记录应包含来源、时间、适用范围和必要的置信标记；用户偏好、事实、推断与临时状态应能区分。摘要会丢失细节，检索会遗漏信息，过时事实可能与新观察冲突，所以应支持更新、撤销和回查原始证据。

外部记忆不是模型参数更新。检索到一段记录只是把它放入本次条件；是否被正确使用仍取决于相关性、表达、位置和可信度。评价记忆模块应同时检查检索命中率、过时记忆干扰和隐私边界，而不只看存储容量。

\section{上下文工程与信息预算}

上下文工程（context engineering）管理系统指令、对话、工具结果、证据、记忆与任务状态如何进入有限窗口。设候选片段价值为 $v_i$、词元成本为 $c_i$，简单选择问题可以写为 $\max\sum_iv_iz_i$，约束 $\sum_ic_iz_i\le B_{\mathrm{ctx}}$、$z_i\in\{0,1\}$。真实片段之间存在冗余和依赖，所以线性背包只是预算视角，不能直接把独立相关性分数当成真实效用。

压缩可以删去无关内容、提取结构化事实或生成摘要。应保留关键约束、证据来源、未完成事项和动作结果；若只保留自然语言结论，后续可能无法区分观察与猜测。上下文窗口还要为输出、工具调用和可能的后续观察预留空间，不能把全部容量填满输入。

过期状态应及时移除或标记。例如旧文件路径、已经失败的假设和已撤销权限若继续存在，可能支配后续决策。上下文管理的目标是在当前任务下保留有效信息，而不是把所有历史都无条件塞入模型。

\section{ReAct、计划执行与自我修正}

ReAct 把推理与动作、观察交替组织，使计划受环境反馈约束\cite{react}。计划--执行（plan-and-execute）先给出较粗子目标，再由执行过程逐步落实；分层规划把大任务拆成不同抽象层级。计划应能在新证据下更新，否则长而详细的初始计划可能成为错误约束。

反思（reflection）与自我修正（self-correction）让模型重新评估先前结果。若没有新证据、外部测试或可靠判据，多次自检可能重复相同错误，甚至把正确答案改坏。因此，修正应围绕具体失败信号，例如测试报错、缺失引用或状态不一致，而非要求模型泛泛“再想一想”。

规划质量应通过任务成功、恢复率、无效动作数、成本和尾部延迟评估。固定工作流是重要基线：如果任务路径本来明确，动态规划可能只增加成本；如果环境变化大、信息逐步揭示，智能体决策才可能体现优势。

\section{多智能体协作与分布式信用}

多智能体（multi-agent）把检索、编码、检查等职责分给多个决策实体。集中式结构由协调者分配任务并合并结果，分散式结构由各实体协商和交接。角色专门化可能让局部上下文更集中，也可能造成信息丢失、重复工作与冲突。

交接（handoff）至少应包含目标、已知事实、已执行动作、剩余约束和可验证产物。共享文件或外部状态需要明确写入责任与合并顺序；两名智能体各自声称完成不代表组合产物一致。投票也不能自动消除相关错误，尤其当模型、提示和数据源高度相似时。

多智能体学习可把联合策略分解为 $\pi(a_1,\ldots,a_M\mid o_1,\ldots,o_M)$ 的受限因子形式，训练时用团队奖励或个体奖励。其他智能体策略变化会使局部环境非平稳，团队回报又使信用分配困难。集中训练、分散执行是一种设计路线，但需要明确训练时可见信息与部署时可见信息。比较多智能体收益必须固定总计算预算，而非只固定每个成员预算。

